\subsection{The actual tensor action and its coherent deformation}
\begin{finding}[\fid{F21}: the specified rank-three tensor lattice is realized]\label{f:r3tensor}
Use the ring $R_0$, algebra $A$, center $C$, roots $z_1,z_2$ and named
generators of \fid{F19}. Let $V=W=R_0^2$ and give
$M=(V\otimes W)^{\otimes3}$ the actual right tensor-monomial action of
GCT~IV, \S6.1, p.~33, and Definition~11.1, equation~(93). Write
$\epsilon_+$ for $P_i=D$, $\epsilon_-$ for $P_i=0$, and
\[
 N(a):\quad P_1=\begin{pmatrix}D&0\\a&0\end{pmatrix},\qquad
 P_2=\begin{pmatrix}0&a\\0&D\end{pmatrix}.
\]
These are right row matrices, as in IV Proposition~11.14,
equation~(110). There is an action-preserving isomorphism over $R_0$,
\begin{equation}\label{eq:r3tensor}
 M\simeq\epsilon_+^{20}\oplus\epsilon_-^4
       \oplus N(\delta)^{16}\oplus N(D-2)^4.
\end{equation}
In particular every displayed change of basis is regular at $v=1$ in
characteristic different from $2,3$; no colliding-root difference is inverted.

Put $L_i=C/(z-z_i)$ and
$S=\Spec R_0\sqcup\Spec R_0\sqcup\Spec C$. For
\[
 E=(\mathcal O,\mathcal O,\mathcal O_C^2),\qquad
 \mathcal F=(\mathcal O^{20},\mathcal O^4,
             \widetilde L_1^{16}\oplus\widetilde L_2^4),
\]
the coherent model realizes this same action as
\[
 M\simeq\Gamma(S,\mathcal Hom(E,\mathcal F)),\qquad A=\Gamma(S,\End E),
\]
with right action by precomposition. It is faithful before specialization
and kills the four-dimensional radical after $v=1$.
\end{finding}
\begin{proof}
Let $L=R_0^2$. The cited ordinary tensor action, with the consistent
normalization $c_i=T_i+v^{-1}$, is
\[
 \mathbf v_{\mathbf k}c_i=
 \begin{cases}
 v\mathbf v_{\mathbf k}+\mathbf v_{\mathbf k s_i},&k_i<k_{i+1},\\
 \delta\mathbf v_{\mathbf k},&k_i=k_{i+1},\\
 v^{-1}\mathbf v_{\mathbf k}+\mathbf v_{\mathbf k s_i},&k_i>k_{i+1}.
 \end{cases}
\]
On $(112,121,211)$ it gives
\[
 c_1=\begin{pmatrix}\delta&0&0\\0&v&1\\0&1&v^{-1}\end{pmatrix},\qquad
 c_2=\begin{pmatrix}v&1&0\\1&v^{-1}&0\\0&0&\delta\end{pmatrix}.
\]
The row-basis matrix
\[
 B_3=\begin{pmatrix}v^2&v&1\\-1-v^{-2}&v&1\\
                         -v^{-1}&-v^{-2}&\delta\end{pmatrix}
\]
has determinant $(D-1)^2$. Multiplying gives
\[
 B_3c_1B_3^{-1}=\begin{pmatrix}\delta&0&0\\0&\delta&0\\0&1&0\end{pmatrix},
 \quad
 B_3c_2B_3^{-1}=\begin{pmatrix}\delta&0&0\\0&0&1\\0&0&\delta\end{pmatrix}.
\]
Thus this block is $\epsilon\oplus U$, with
\[
 c_1|_U=\begin{pmatrix}\delta&0\\1&0\end{pmatrix},\qquad
 c_2|_U=\begin{pmatrix}0&1\\0&\delta\end{pmatrix}.
\]
On $(122,212,221)$ the original $c_1,c_2$ matrices interchange; applying
$B_3$ and swapping its last two rows again gives $\epsilon\oplus U$.
The weight spaces $111$ and $222$ are each $\epsilon$. Since the determinants
are units, this proves $L^{\otimes3}=\epsilon^4\oplus U^2$ over $R_0$,
not just over its fraction field.

By $P_i=c_i\otimes c_i+(c_i-\delta)\otimes(c_i-\delta)$,
$\epsilon\otimes\epsilon=\epsilon_+$, and each mixed tensor is
$N(\delta)$. On $U\otimes U$ in the order $(11,12,21,22)$,
\[
 P_1=\begin{pmatrix}D&0&0&0\\\delta&0&0&0\\\delta&0&0&0\\
                         2&-\delta&-\delta&D\end{pmatrix},\quad
 P_2=\begin{pmatrix}D&-\delta&-\delta&2\\0&0&0&\delta\\
                         0&0&0&\delta\\0&0&0&D\end{pmatrix}.
\]
Take
\[
 B_4=\begin{pmatrix}
 \delta&-1&-1&\delta\\0&-1&1&0\\
 2-D&-\delta&-\delta&D\\-D&\delta&\delta&D-2
 \end{pmatrix},\qquad \det B_4=8(D-1)^2.
\]
Direct multiplication yields
\[
 B_4P_1B_4^{-1}=\begin{pmatrix}D&0&0&0\\0&0&0&0\\0&0&D&0\\0&0&D-2&0\end{pmatrix},
\quad
 B_4P_2B_4^{-1}=\begin{pmatrix}D&0&0&0\\0&0&0&0\\0&0&0&D-2\\0&0&0&D\end{pmatrix}.
\]
This is $\epsilon_+\oplus\epsilon_-\oplus N(D-2)$, again over $R_0$.
In $(\epsilon^4\oplus U^2)\otimes(\epsilon^4\oplus U^2)$ the unmixed
trivial contribution has multiplicity $16$, the two mixed contributions
have total multiplicity $16$, and $U\otimes U$ occurs four times.
This proves \eqref{eq:r3tensor}, whose rank is $20+4+32+8=64$.

It remains to match the named action to the center model, rather than
merely match dimensions. For $a=\delta,D-2$, put $q_a=a/D$, a unit, and
\[
 S_a=\begin{pmatrix}1&0\\-q_a^2&q_a\end{pmatrix}.
\]
The matrices of $S_a(P_1/D)S_a^{-1}$ and $S_a(P_2/D)S_a^{-1}$ are
respectively
\[
 \begin{pmatrix}1&0\\0&0\end{pmatrix},\qquad
 \begin{pmatrix}1-z_i&1\\z_i(1-z_i)&z_i\end{pmatrix},
 \quad z_i=1-q_a^2.
\]
They are exactly the standard row module for $M_2(C)$ evaluated on branch
$i$. The right-module formula with $\mathcal Hom$ follows. The two branch
evaluations jointly embed $C$ into $R_0^2$, so their modules have joint
annihilator zero; both scalar blocks also occur. At $v=1$ the branch
modules become the same residue module of $k[\varepsilon]/\varepsilon^2$,
and the central nilpotent acts by zero. Its matrix ideal is the radical
of \fid{F19}, proving the last assertion.
\end{proof}
The generic multiplicities themselves are already in IV \S12.3; the
additional assertion here is the displayed regular action isomorphism.
The coherent sheaf $\mathcal F$ is not locally free on the collision
component: its two generic branch ranks are $16$ and $4$. This result
does not identify a simultaneous canonical left quantum-group action
or VIII's based tensor lattice with the monomial lattice.

\begin{finding}[\fid{F22}: constant nearby-cycle purity does not detect collision order]\label{f:collision}
Let $R=k[[t]]$ and, for $m\ge1$, set
\[
 C_m=\{(u,w)\in R^2:u\equiv w\pmod{t^m}\}
     \simeq R[x]/(x(x-t^m)).
\]
For $m\ge n\ge1$, $C_m\subset C_n$ induces a finite universal
homeomorphism $\Spec C_n\to\Spec C_m$ over $\Spec R$, although
\[
 \operatorname{length}_R(R^2/C_m)=m,\qquad
 \operatorname{cond}(C_m,R^2)=t^mR^2.
\]
Their corresponding ordinary \'etale coefficient categories and
cohomology are equivalent. With $k=\F_q$ and unshifted constant
$\Q_\ell$ coefficients, the unique special-point nearby and vanishing
cycles are
\[
 R\Psi\Q_\ell=\Q_\ell^2[0],\qquad
 R\Phi\Q_\ell=\Q_\ell[0],
\]
with trivial inertia and Frobenius. These constant data, including every
finite-extension trace and tensor power, are independent of $m$.
The completed center of \fid{F19} in good characteristic is $C_2$.
\end{finding}
\begin{proof}
Map $x$ to $(0,t^m)$ to obtain the displayed description. For $m\ge n$,
$C_n$ is a finite $C_m$-module since it is already generated by two
elements over the subring $R$. After inverting $t$, the map is the
identity on $k((t))^2$. Each special fiber has one point, with residue
field $k$, and the map induces the identity on that field. It is
therefore integral, surjective and radicial, hence a universal
homeomorphism. The difference map
$(u,w)\mapsto w-u\bmod t^m$ has kernel $C_m$ and gives
$R^2/C_m\simeq R/(t^m)$. The integral closure is $R^2$; a pair multiplies
every element of $R^2$ back into $C_m$ precisely when both coordinates
are divisible by $t^m$. This proves the coherent formulas.

Topological invariance of the \'etale site gives the claimed equivalence,
compatible with the base trait; precise locators are Stacks
Theorem~59.45.2, tag~04DZ, and Proposition~59.45.4, tag~03SI,
the former citing EGA~IV, 18.1.2 \cite{StacksNil}.
This is an equivalence of coefficient-sheaf topoi, not an identification
of coherent structure sheaves as ringed objects.

After strict henselization at the special point, the geometric generic
fiber is two points. Indeed the finite local algebra over the strictly
henselian trait is strictly henselian, and its two branches are already
split over the fraction field. The defining nearby-cycle stalk is
therefore the cohomology of two geometric points: $\Q_\ell^2$ in degree
zero, with trivial inertia and Frobenius. The specialization map from
$i^*\Q_\ell=\Q_\ell[0]$ is diagonal. The defining triangle
\[
 i^*\Q_\ell\longrightarrow R\Psi\Q_\ell
 \longrightarrow R\Phi\Q_\ell\longrightarrow i^*\Q_\ell[1]
\]
thus gives the stated cokernel in degree zero and no other groups.
This also follows by finite base change, Stacks Lemma~59.91.9,
tag~0A4E. Every object and map just described is independent of $m$.
The traces of tensor powers $r$ of $R\Psi$ and $R\Phi$ are respectively
$2^r$ and $1$ over every finite extension.

Finally \eqref{eq:r3collision} says that the actual root difference is
a unit times $(v-1)^2$. Translate $z$ by a root and rescale the translated
coordinate by that unit after completion. The relation becomes
$x(x-t^2)=0$, where $t=v-1$. This identifies the center with $C_2$.
\end{proof}
This is an information limit for the specified constant coefficient
families. A nilpotent algebra action supplied separately, or an
$m$-dependent coefficient construction, is additional information and
is not ruled out. A sign criterion applied only to these common traces
can establish their correct weight-zero purity but cannot determine
$m$. The next finding exhibits coherent information that does retain it.

\begin{finding}[\fid{F25}: the commutant defect, coherent extensions, and a regular basis]\label{f:r3deformation}
Use \fid{F21} over $R=\Q[v,v^{-1}]_{(v-1)}$ and write
$\Delta=z_1-z_2$. Then
\[
 \End_A(M)=M_{20}(R)\times M_4(R)\times M_{16}(R)\times M_4(R).
\]
This is free of rank $688$. Its base-change map to the full special
commutant
\[
 \End_{A_1}(M_1)=M_{20}(k)\times M_4(k)\times M_{20}(k)
\]
is injective with cokernel of dimension $128$; the latter commutant has
dimension $816$. Moreover
\[
 \operatorname{Ext}^1_C(L_1,L_2)
 =\operatorname{Ext}^1_C(L_2,L_1)=R/(\Delta),\qquad
 \operatorname{Ext}^1_C(L_i,L_i)=0,
\]
and $\operatorname{Ext}^1_A(M,M)=(R/(\Delta))^{128}$, of $R$-length $256$.

There is also a regular center basis $1,y$, with
\begin{equation}\label{eq:r3basis}
 y=D^2(z_1-z),\qquad
 y^2=\lambda y,\qquad
 \lambda=(v-v^{-1})^2(v^2+1+v^{-2}).
\end{equation}
Together with the two scalar-block units, $E_{ij},yE_{ij}$ form a basis
of the actual rank-ten algebra over $R_0$, with structure constants
among $0,1,\lambda$. This elementary basis is not asserted to be VIII's
designated canonical basis.
\end{finding}
\begin{proof}
Central blocks separate the two scalar sectors and the matrix sector.
Morita equivalence identifies the last sector's multiplicity module
with $L_1^{16}\oplus L_2^4$. A homomorphism $L_1\to L_2$ is determined
by $r\in R$ satisfying $(z_2-z_1)r=0$. Since $\Delta$ is a nonzero
element of the domain $R$, this Hom group is zero, as is the reverse
one; $\End_C(L_i)=R$. This gives the asserted commutant. At $v=1$ the
two branch modules become the same simple residue module. Their
multiplicity spaces merge into $k^{20}$. The map on these factors is
the block-diagonal embedding $M_{16}(k)\times M_4(k)\hookrightarrow
M_{20}(k)$; the missing off-diagonal blocks each have dimension $16\cdot4$.

For the extension calculation put $f_i=z-z_i$. The following is an
exact free resolution of $L_1$:
\[
 \cdots\xrightarrow{f_1}C\xrightarrow{f_2}C
             \xrightarrow{f_1}C\longrightarrow L_1\longrightarrow0.
\]
To check exactness, lift to the domain $R[z]$ and cancel $f_1$ or $f_2$
from divisibility by $f_1f_2$. Applying $\Hom_C(-,L_2)$ gives
$R\xrightarrow{-\Delta}R\xrightarrow0R\xrightarrow{-\Delta}R\to\cdots$,
so degree one is $R/(\Delta)$. Exchanging indices gives the reverse
group. Applying $\Hom_C(-,L_1)$ instead begins
$R\xrightarrow0R\xrightarrow\Delta R$ and has zero degree one.
Morita equivalence and the multiplicities in \eqref{eq:r3tensor} give
$16\cdot4+4\cdot16=128$ copies; scalar blocks contribute no extensions
of their free $R$-modules. Since $\Delta$ is a unit times $(v-1)^2$,
each crossed quotient has length two. For the general $C_m$ the same
resolution gives crossed extension group $R/(t^m)$, retaining precisely
the collision order omitted by the constant \'etale data.

Finally $D$ is a unit and the affine change $z\mapsto y$ is invertible.
The two evaluations of $y$ are $0$ and $D^2(z_1-z_2)$. Their product
relation yields $y^2=D^2\Delta y$. Since
\[
 D^2\Delta=(D-1)(D-4),\quad
 D-4=(v-v^{-1})^2,\quad D-1=v^2+1+v^{-2},
\]
this is \eqref{eq:r3basis}. Matrix-unit multiplication proves the
structure-constant assertion. At $v=1$, $\lambda=0$ and the four
$yE_{ij}$ span the radical.
\end{proof}
IV \S12.3 already records the numbers $688$ and $816$. Here the free
rank-$688$ commutant family is kept distinct from the full
special-fiber commutant, and their comparison map is explicit.

The scalar in \eqref{eq:r3basis} has a precise ordinary geometric trace:
with $Q=v^2$,
\[
 \lambda=Q^{-2}(Q-1)^2(1+Q+Q^2)
       =Q^{-2}\#(\mathbb G_m^2\times\PP^2)(\F_Q).
\]
It is the compact-support alternating trace with Tate twist $(2)$.
The remaining factor $v^{-2}(1+v^2+v^4)$ is the centered Tate polynomial
of $\PP^2$. This mixed-cohomology trace interpretation of one scalar
does not supply functorial multiplication or associativity constraints,
and it does not assert VIII's consecutive-$v$ unimodality convention.
Indeed $\lambda=v^4-v^2-v^{-2}+v^{-4}$ has negative Laurent coefficients.

